\section{Unified Research}
\label{sec:research}

\subsection{Research as on-chain artifacts}

Research outputs that influence the model---architecture variants, hyperparameter choices, ablation results, benchmark scores---are stored as first-class on-chain artifacts. The \texttt{RegisterArtifact} instruction commits an artifact of one of four kinds:

\begin{longtable}{p{3cm}p{4cm}p{7cm}}
\toprule
\textbf{Kind} & \textbf{Hashed object} & \textbf{Use} \\
\midrule
\endhead
\texttt{architecture} & model definition file & defines the structure trained in subsequent epochs \\
\texttt{hyperparams} & training config (yaml/json) & locks the schedule used for the run \\
\texttt{ablation} & experiment log + result tuple & records what was tried and the outcome \\
\texttt{benchmark} & dataset commitment + scores & committed evaluation result \\
\bottomrule
\end{longtable}

\subsection{Research artifact registry contract}

The chain ships with a \texttt{ResearchRegistry} precompile-backed contract whose state-changing methods are restricted to addresses on the model's contributor allowlist (a list managed by the model's governance). Reads are unrestricted.

\begin{lstlisting}[language=Solidity,caption=Registry interface (excerpt)]
interface IResearchRegistry {
    function commitArchitecture(bytes32 hash, bytes calldata meta) external;
    function commitHyperparams(uint256 epoch, bytes32 hash) external;
    function commitAblation(bytes32 parentArch, bytes32 hash, int256 deltaScore) external;
    function commitBenchmark(bytes32 dataset, bytes32 scoresHash) external;
    function get(bytes32 artifactId) external view returns (Artifact memory);
}
\end{lstlisting}

\subsection{Reproducibility from artifacts}

Because the artifact registry, the data tree, and the receipts tree are all on the same chain, reproducing any reported result is mechanical:
\begin{enumerate}[nosep]
\item Look up the architecture artifact and hyperparameters for the target epoch.
\item Walk the receipts tree to enumerate the gradient transactions that built the epoch's checkpoint.
\item Replay deterministically per \S\ref{sec:training}.
\item Look up the benchmark artifact and re-run on the committed dataset.
\end{enumerate}

The chain itself does not run benchmarks (they require model inference, which the \Hanzo{} AI Chain serves), but it provides a tamper-evident bibliography for any reported number.

\subsection{Citation graph}

Ablations form a directed graph rooted at architecture artifacts. The graph is queryable via standard subgraph indexers; the chain emits typed events on every commit. This produces a transparent research history: any claim of the form ``architecture X improves Y by $\delta$'' has a chain transaction that records the parent artifact, the child artifact, and the delta.

\subsection{Research credit attribution}

Research artifacts contribute to the attribution tree the same way training shards do, but with a configurable multiplier. The default policy:
\begin{equation}
w_{\text{author} \to \text{epoch}} \mathrel{+}=
\begin{cases}
\alpha_{\text{arch}} \cdot \mathbf{1}[\text{arch adopted}] & \text{if architecture} \\
\alpha_{\text{hp}} \cdot \mathbf{1}[\text{hp adopted}] & \text{if hyperparams} \\
\alpha_{\text{abl}} \cdot \max(0, \delta) & \text{if ablation} \\
\alpha_{\text{bench}} \cdot \log(1 + \text{citations}) & \text{if benchmark}
\end{cases}
\end{equation}

with $\alpha_{\text{arch}} = 1.0$, $\alpha_{\text{hp}} = 0.1$, $\alpha_{\text{abl}} = 0.5$, $\alpha_{\text{bench}} = 0.2$ as defaults. The values are amendable by chain governance.
